% Generated draft: requires mathematical and source review.
\documentclass[11pt]{article}
\usepackage[margin=25mm]{geometry}
\usepackage{fontspec}
\setmainfont{Latin Modern Roman}
\usepackage{amsmath,amssymb,mathtools,unicode-math}
\setmathfont{Latin Modern Math}
\usepackage{enumitem,needspace,xurl,hyperref}
\setlength{\parindent}{0pt}
\setlength{\parskip}{5pt}
\setlength{\emergencystretch}{3em}
\newcommand{\N}{\mathbb{N}}
\newcommand{\Z}{\mathbb{Z}}
\newcommand{\Q}{\mathbb{Q}}
\newcommand{\R}{\mathbb{R}}
\newcommand{\C}{\mathbb{C}}
\newcommand{\F}{\mathbb{F}}
\newcommand{\A}{\mathbb{A}}
\newcommand{\PP}{\mathbb{P}}
\newcommand{\E}{\mathbb{E}}
\newcommand{\eps}{\varepsilon}
\newcommand{\dd}{\,\mathrm{d}}
\begin{document}
\section*{UnsolvedMath ID: 700426}
\textbf{Source problem code:} catalogue entry
\subsection{Definitions and mathematical statement}
Let \(n\geq 2\) be an integer and let \(q\) be a prime power. Write \(V=\F_q^n\) and
\[
\operatorname{SL}_n(q)=\{g\in\operatorname{GL}_n(\F_q):\det g=1\}.
\]
A subgroup \(H\leq\operatorname{SL}_n(q)\) is irreducible on \(V\) if its only invariant \(\F_q\)-linear subspaces are \(0\) and \(V\). For this draft, call \(H\) block-monomial if there is a decomposition
\[
V=V_1\oplus\cdots\oplus V_r,\qquad r\geq2,\qquad \dim_{\F_q}V_i=b\geq1,
\]
with \(rb=n\), such that every element of \(H\) permutes the summands. Equivalently, in a basis adapted to this decomposition, every element has exactly one nonzero, invertible \(b\times b\) block in each block row and column.

\textbf{Question (V. M. Levchuk).} With these interpretations, is it true that
\[
\begin{gathered}
\forall n\in\Z,\ n\geq2,\quad \exists Q(n)\in\N\quad
\forall q\text{ a prime power with }q\geq Q(n),\\
\exists a_1,\ldots,a_n\in\F_q^{\times},\qquad
\prod_{i=1}^{n}a_i=1,
\end{gathered}
\]
such that, for \(d=\operatorname{diag}(a_1,\ldots,a_n)\), every subgroup \(H\) satisfies
\[
\bigl(H<\operatorname{SL}_n(q),\ d\in H,\ H\text{ irreducible on }V\bigr)
\quad\Longrightarrow\quad H\text{ is block-monomial}?
\]
Here \(<\) means proper subgroup, and diagonal means diagonal in the standard basis of \(V\).

\textbf{Unresolved qualifications.} The source does not define ``irreducible'' or ``block-monomial.'' The definitions above use irreducibility over \(\F_q\) and a block decomposition that may depend on \(H\); they require confirmation against the original source. The quantifiers interpret ``sufficiently large \(q\)'' for each fixed \(n\), allowing the threshold to depend on \(n\); the source does not explicitly state this dependence or the range of \(n\).

The supplied source reports an affirmative answer for \(n=2,3\), attributed to V. M. Levchuk's cited 1973 contribution. This is a supplied historical assertion, not an independently verified result here. Its additional suggestion of similar questions for other Chevalley groups does not specify groups, representations, or an analogue of the exceptional subgroups, so it supplies no further precise statement.
\subsection{Short English statement}
For each fixed dimension \(n\geq2\), ask whether every sufficiently large finite field admits a determinant-one diagonal matrix whose proper irreducible overgroups are all block-monomial. The meanings of irreducibility and block-monomial, and the dependence of the field-size threshold on \(n\), require source confirmation. The source reports affirmative answers in dimensions \(2\) and \(3\).
\subsection{Sources}
\begin{itemize}
\item Local UnsolvedMath registry, permanent ID \(700426\), source code \texttt{catalogue entry}, title \emph{Kourovka Notebook Problem 7.27}; supplied revision \texttt{d371ee917ac9647fbb06cd6fc529deec6f3dda82}. The registry labels the problem \emph{open}; this supplied classification is not independently verified here.
\item \emph{The Kourovka Notebook (Unsolved Problems in Group Theory)}, problem 7.27, proposed by V. M. Levchuk, 7th issue (1980). The supplied record identifies the 21st issue edition (2026), \texttt{arXiv:1401.0300v48}; the collection metadata names E. I. Khukhro and V. D. Mazurov as editors.
\item V. M. Levchuk, contribution in \emph{Some problems of the theory of groups and rings}, Inst. of Physics SO AN SSSR, Krasnoyarsk, 1973 (Russian). The supplied statement cites this reference for the affirmative cases \(n=2,3\), without a contribution title, page range, or theorem number.
\end{itemize}
\end{document}
